\section{Availability and Storage Scalability of Archive Layer}

\texttt{get} shards for versions that have not been garbage collected always succeeds because out of the $k+2f$ queried nodes, at least $k+f$ are correct.
Since any $k$ correct nodes holding any $k$ pieces of the stripe (state or parity shards) can reconstruct the original $k$ data shards, the queried stripes can always be reconstructed.

The garbage collection process requires a bump quorum of only $2f+1$ participating nodes.
As mentioned in \autoref{sec:approach}, for each stripe, the $k+2f$ nodes responsible for serving it must overlap with this $2f+1$ quorum by at least $k$ correct nodes.
These $k$ correct nodes will possess enough state and parity shards to reconstruct the entire stripe.

As the result, \texttt{get} every shard of the \emph{highest} version certified by a bump quorum will succeed.
Simultaneously, the bump quorum certificate also indicates the correct Merkle root (view of the state) for that version.
So, if the incoming query is for a shard versioned older than the last bump quorum certificate, the nodes can respond with the certificate instead of the shard.
This causes the corresponding \texttt{get} to return null, but in a following \texttt{skip} call, the querying node can fast-forward to the certified version and set its view of the state to the certified Merkle root.
Therefore, the nodes can safely discard shards older than the last certified version.

\paragraph{Redundancy Analysis.}
The redundancy introduced by the EC$(k+2f,k)$ erasure code is $\frac{k+2f}{k}$.
When $k=1$, the redundancy becomes $2f+1$, which is equivalent to the plain replication design I1.
By setting $k=pf$, where $p$ is a constant in the range $(0,\frac{f+1}{f}]$, the redundancy becomes $\frac{pf+2f}{pf} = \frac{p+2}{p}$. This redundancy factor remains constant regardless of the state size and total storage capacity, thus achieving the desired storage scalability property.

However, this also implies that as $f$ increases (\ie, more nodes join the system), $k$ should ideally increase as well to maintain a target redundancy ratio.
This necessitates dynamically adjusting the number of shards within a stripe.
% Fortunately, advancements in convertible codes~\cite{maturana2020access,maturana2022convertible,maturana2023bandwidth} and StripeMerge~\cite{yao2021stripemerge} offer a solution.
Fortunately, we can employ a specialized erasure code~\cite{maturana2020access,maturana2022convertible,maturana2023bandwidth,yao2021stripemerge} that allows merging multiple stripes into a single larger stripe (with a proportionally increased $k$) through efficient computation by combining the parity shards of the original stripes.
This approach mitigates the significant overhead of massive transcoding operations.

It's important to note that even if $k$ is not adjusted immediately with an increasing $f$, the safety of the system is not compromised; we would simply incur a higher redundancy overhead than necessary.
Conversely, when $f$ decreases, $k$ must be reduced, and the stripes need to be reconstructed promptly.

Nevertheless, when the total number of nodes $n$ increases, the potential number of faults $f$ also increases, but an immediate increase in $k$ is not strictly required.
Conversely, when $n$ decreases, $f$ does not necessarily need to be decreased immediately (except for potential liveness considerations), and the administrator can choose an opportune time to decrease $f$ and consequently $k$.
In summary, $k$ does not require immediate adjustment regardless of whether $n$ is increasing or decreasing.
In future work, we will design a comprehensive reconfiguration protocol leveraging these techniques to achieve low-overhead scaling up/down.
\lijl{I assume this should be part of the full paper?}\gd{Yes}